% =====================================================================
\section{Choques oblicuos: parámetro de impacto}
\label{sec:oblicuas}
% =====================================================================

\begin{definicion}[Parámetro de impacto]
\label{def:b}
Es la distancia $b$ entre la recta que sigue el centro del proyectil y la paralela que
pasa por el centro del blanco. Para esferas de radios $R_{1}$, $R_{2}$ hay choque si
$b<d\defeq R_{1}+R_{2}$, y en el contacto la línea de centros forma con la dirección
incidente el ángulo
\begin{equation}
  \resultado{\sin\varphi=\frac bd} .
\end{equation}
Se toma $\uv{n}$ a lo largo de la línea de centros (de $1$ hacia $2$) y $\uv{t}\perp\uv{n}$ en el
plano del choque (\cref{fig:oblicuo}).
\end{definicion}

\begin{proposicion}[Esferas lisas]
\label{prop:lisas}
Si las superficies son lisas, la percusión es normal ($\bm{\pi}\parallel\uv{n}$) y
\begin{enumerate}[label=\textup{(\alph*)},leftmargin=2em]
  \item las componentes tangenciales no cambian: $v_{it}'=v_{it}$;
  \item las normales siguen el choque frontal \eqref{eq:frontal}: $v_{in}'=(1+e)V_{n}-e\,v_{in}$,
        con $V_{n}=\Vcm\cdot\uv{n}$;
  \item $Q=\tfrac12\,\mu\,(1-e^{2})\,(\vrel\cdot\uv{n})^{2}$.
\end{enumerate}
\end{proposicion}

\begin{proof}
$m_{1}\Delta\vb{v}_{1}=-\pi\,\uv{n}$ y $m_{2}\Delta\vb{v}_{2}=\pi\,\uv{n}$: proyectando sobre $\uv{t}$, $\Delta v_{it}=0$.
Sobre $\uv{n}$ el problema es unidimensional y $e$ se define con las componentes normales,
luego vale el \cref{teo:frontal}. En $K=\tfrac12MV\cm^{2}+\tfrac12\mu(v_{12,n}^{2}+v_{12,t}^{2})$ sólo
cambia el término normal, que se multiplica por $e^{2}$.
\end{proof}

\begin{clave}{Método · choque oblicuo}
\begin{enumerate}[leftmargin=1.4em]
  \item Geometría: $\sin\varphi=b/d$; base $\{\uv{n},\uv{t}\}$ ligada a la línea de centros.
  \item Descomponer las velocidades en $\uv{n}$ y $\uv{t}$.
  \item Normal: choque frontal con $e$. Tangencial: se conserva.
  \item Recomponer en la base del \lab{} y calcular ángulos con $\tan=v_{y}/v_{x}$.
\end{enumerate}
\end{clave}

% ---------------------------------------------------------------------
\begin{figure}[ht]
\centering
\begin{tikzpicture}[scale=1.5]
  \def\ph{35}\def\d{1.2}\def\Ra{0.5}\def\Rb{0.7}
  \coordinate (C2) at (0,0);
  \coordinate (C1) at ({-\d*cos(\ph)},{\d*sin(\ph)});
  \coordinate (P) at ($(C1)!\Ra/\d!(C2)$); % punto de contacto
  % rectas de referencia
  \draw[aux] (-3.6,0) -- (2.6,0);
  \draw[aux] (-3.6,{\d*sin(\ph)}) -- ($(C1)+(1.55,0)$);
  \draw[tenue!80,line width=0.4pt] ($(C2)!-0.35!(C1)$) -- ($(C2)!1.3!(C1)$);
  % esferas
  \fill[rosa,opacity=0.20] (C2) circle (\Rb); \draw[rosa,line width=0.7pt] (C2) circle (\Rb);
  \fill[cian,opacity=0.20] (C1) circle (\Ra); \draw[cian,line width=0.7pt] (C1) circle (\Ra);
  \node[punto,inner sep=0.9pt] at (C2) {}; \node[punto,inner sep=0.9pt] at (C1) {};
  \node[etiq,cian,anchor=south,inner sep=2pt] at ($(C1)+(90:\Ra)$) {$m_{1}$};
  \node[etiq,rosa,anchor=north east,inner sep=1pt] at ($(C2)+(220:\Rb+0.08)$) {$m_{2}$};
  % parámetro de impacto
  \draw[tenue,line width=0.45pt,{Stealth[length=3.5pt]}-{Stealth[length=3.5pt]}] (-3.25,0) -- (-3.25,{\d*sin(\ph)})
    node[midway,anchor=east,etiq,inner sep=3pt] {$b$};
  % velocidad incidente
  \draw[vec,cian] (-2.85,{\d*sin(\ph)}) -- (-1.75,{\d*sin(\ph)}) node[midway,above,etiq,cian] {$\vb{v}_{1}$};
  % base n, t en el contacto
  \draw[vec,verde,line width=0.7pt] (P) -- ++({-\ph}:0.42);
  \node[etiq,verde,anchor=north east,inner sep=1pt] at ($(P)+({-\ph}:0.36)$) {$\uv{n}$};
  \draw[vec,verde,line width=0.7pt] (P) -- ++({90-\ph}:0.42) node[anchor=south west,etiq,verde,inner sep=0.5pt] {$\uv{t}$};
  % velocidades finales
  \draw[vec,rosa] ($(C2)+({-\ph}:\Rb)$) -- ($(C2)+({-\ph}:2.0)$) node[anchor=north west,etiq,rosa,inner sep=2pt] {$\vb{v}_{2}'$};
  \draw[ambar,line width=0.55pt] (1.2,0) arc (0:{-\ph}:1.2);
  \node[etiq,ambar,anchor=west,inner sep=2pt] at ({-\ph/2}:1.25) {$\beta=\varphi$};
  \draw[vec,cian] ($(C1)+(39:\Ra)$) -- ($(C1)+(39:2.0)$) node[anchor=south west,etiq,cian,inner sep=2pt] {$\vb{v}_{1}'$};
  \draw[ambar,line width=0.55pt] ($(C1)+(1.35,0)$) arc (0:39:1.35);
  \node[etiq,ambar] at ($(C1)+(19.5:1.58)$) {$\alpha$};
\end{tikzpicture}
\caption{Choque oblicuo de esferas lisas en el instante del contacto. La línea de centros
forma el ángulo $\varphi$ con la dirección incidente ($\sin\varphi=b/d$); la percusión va según
$\uv{n}$, luego el blanco sale por la línea de centros, con $\beta=\varphi$ (dibujado con $m_{1}=m_{2}$, $e=0{,}6$).}
\label{fig:oblicuo}
\end{figure}
% ---------------------------------------------------------------------

\begin{problema}{2}[Ángulo de difusión]
\label{prob:2}
Obtener el ángulo de difusión en una colisión oblicua de $m_{1}$ y $m_{2}$ con parámetro de
impacto $b$, si $m_{2}$ está en reposo y $m_{1}$ incide con $\vb{v}_{1}$. Coeficiente de restitución $e$.
\end{problema}

\begin{resolucion}
Esferas lisas, $d=R_{1}+R_{2}$, $\vb{v}_{1}=v_{1}\ii$, $M=m_{1}+m_{2}$.
\begin{enumerate}
  \item \emph{Base}: $\sin\varphi=b/d$, $\cos\varphi=\sqrt{d^{2}-b^{2}}/d$;
        $\uv{n}=(\cos\varphi,-\sin\varphi)$, $\uv{t}=(\sin\varphi,\cos\varphi)$.
  \item \emph{Componentes iniciales}: $v_{1n}=v_{1}\cos\varphi$, $v_{1t}=v_{1}\sin\varphi$; $\vb{v}_{2}=\vb{0}$.
  \item \emph{Normal} (\cref{teo:frontal} con $v_{2n}=0$):
        \[
          v_{1n}'=\frac{m_{1}-e\,m_{2}}{M}\,v_{1}\cos\varphi,\qquad
          v_{2n}'=\frac{(1+e)\,m_{1}}{M}\,v_{1}\cos\varphi .
        \]
        \emph{Tangencial}: $v_{1t}'=v_{1}\sin\varphi$, $v_{2t}'=0$.
  \item \emph{Blanco}: $\vb{v}_{2}'=v_{2n}'\uv{n}$, luego sale por la línea de centros:
        $\resultado{\beta=\varphi=\arcsin(b/d)}$.
  \item \emph{Proyectil}: $\vb{v}_{1}'=v_{1n}'\uv{n}+v_{1t}'\uv{t}$, es decir,
        \begin{align*}
          v_{1x}'&=v_{1n}'\cos\varphi+v_{1}\sin^{2}\varphi
                 =\frac{v_{1}}{M}\bigl[m_{1}+m_{2}\sin^{2}\varphi-e\,m_{2}\cos^{2}\varphi\bigr],\\
          v_{1y}'&=-v_{1n}'\sin\varphi+v_{1}\sin\varphi\cos\varphi
                 =\frac{v_{1}}{M}\,(1+e)\,m_{2}\sin\varphi\cos\varphi .
        \end{align*}
  \item \emph{Ángulo}: $\tan\alpha=v_{1y}'/v_{1x}'$:
        \[
          \resultado{\tan\alpha=\frac{(1+e)\,m_{2}\sin\varphi\cos\varphi}{m_{1}+m_{2}\sin^{2}\varphi-e\,m_{2}\cos^{2}\varphi}
          =\frac{(1+e)\,m_{2}\,b\sqrt{d^{2}-b^{2}}}{m_{1}d^{2}+m_{2}b^{2}-e\,m_{2}(d^{2}-b^{2})}} .
        \]
        Si el denominador es negativo, $\alpha>\pi/2$ (el proyectil retrocede).
  \item \emph{Comprobaciones}: $e=1$, $m_{1}=m_{2}$: $\tan\alpha=\cot\varphi$, luego $\alpha+\beta=\pi/2$ (\cref{cor:casos});
        $b\to d$ (roce): $\alpha\to0$; $b=0$: choque frontal, $\alpha=0$ o $\pi$.\qedhere
\end{enumerate}
\end{resolucion}

\begin{problema}{5}[Energía liberada]
\label{prob:5}
Obtener la fracción de energía liberada en la colisión oblicua de $m_{1}$ y $m_{2}$, con $m_{2}$
inicialmente en reposo, parámetro de impacto $b$ y coeficiente de restitución $e$.
\end{problema}

\begin{resolucion}
\begin{enumerate}
  \item La energía tangencial no cambia (\cref{prop:lisas}); sólo se pierde energía del
        movimiento \emph{normal relativo}, con $v_{12,n}=v_{1}\cos\varphi$:
        \[
          Q=\tfrac12\,\mu\,(1-e^{2})\,v_{1}^{2}\cos^{2}\varphi .
        \]
  \item Dividiendo por $T_{1}=\tfrac12m_{1}v_{1}^{2}$, con $\mu/m_{1}=m_{2}/M$ y $\cos^{2}\varphi=1-b^{2}/d^{2}$:
        \[
          \resultado{\frac{Q}{T_{1}}=\frac{m_{2}\,(1-e^{2})}{m_{1}+m_{2}}\Bigl(1-\frac{b^{2}}{d^{2}}\Bigr)} .
        \]
  \item \emph{Comprobación directa}: con los resultados del \cref{prob:2},
        $T_{1}-\tfrac12m_{1}(v_{1n}'^{2}+v_{1}^{2}\sin^{2}\varphi)-\tfrac12m_{2}v_{2n}'^{2}
        =\tfrac12m_{1}v_{1}^{2}\cos^{2}\varphi-\tfrac12m_{1}v_{1n}'^{2}-\tfrac12m_{2}v_{2n}'^{2}$, que es la $Q$
        de un choque frontal con velocidad $v_{1}\cos\varphi$.
  \item \emph{Límites}: máxima en el frontal ($b=0$), nula en el roce ($b\to d$) y nula si $e=1$.\qedhere
\end{enumerate}
\end{resolucion}

% ---------------------------------------------------------------------
\begin{figure}[ht]
\centering
\begin{tikzpicture}
\begin{groupplot}[group style={group size=2 by 1,horizontal sep=16mm},
  oscuro,width=7.2cm,height=5.8cm,xmin=0,xmax=1,xlabel={$b/d$}]
  \nextgroupplot[ymin=0,ymax=95,ytick={0,15,30,45,60,75,90},ylabel={ángulo (grados)},
    title={(a) Difusión y retroceso},legend to name=leyob,legend columns=4,
    legend style={/tikz/every even column/.append style={column sep=8pt}}]
    \addplot[cian,domain=0:0.999,samples=150] {90-asin(x)};
    \addlegendentry{$\alpha$, $e=1$}
    \addplot[ambar,domain=0:0.999,samples=150] {atan2(1.6*x*sqrt(1-x^2),1+x^2-0.6*(1-x^2))};
    \addlegendentry{$\alpha$, $e=0{,}6$}
    \addplot[violeta,domain=0:0.999,samples=150] {atan2(x*sqrt(1-x^2),1+x^2)};
    \addlegendentry{$\alpha$, $e=0$}
    \addplot[rosa,dashed,domain=0:1,samples=150] {asin(x)};
    \addlegendentry{$\beta$}
  \nextgroupplot[ymin=0,ymax=0.55,ylabel={$Q/T_{1}$},
    title={(b) Energía liberada},legend pos=north east]
    \addplot[violeta,domain=0:1,samples=100] {0.5*(1-x^2)};
    \addlegendentry{$e=0$}
    \addplot[ambar,domain=0:1,samples=100] {0.5*0.64*(1-x^2)};
    \addlegendentry{$e=0{,}6$}
    \addplot[cian,domain=0:1,samples=100] {0};
    \addlegendentry{$e=1$}
\end{groupplot}
\node[anchor=north] at ($(group c1r1.south east)!0.5!(group c2r1.south west)+(0,-10mm)$) {\pgfplotslegendfromname{leyob}};
\end{tikzpicture}
\caption{Choque oblicuo con $m_{1}=m_{2}$ (\cref{prob:2,prob:5}). (a) El blanco sale siempre
con $\beta=\arcsin(b/d)$; el proyectil se desvía menos cuanto más inelástico es el choque.
(b) $Q/T_{1}=\tfrac12(1-e^{2})(1-b^{2}/d^{2})$.}
\label{fig:oblicuo-curvas}
\end{figure}
% ---------------------------------------------------------------------
